\section{Sharp spectral accuracy for replacing the obstacle matrix}
\label{sec:op3-spectral-envelope-stability}

Spectral approximation can accelerate an accurate computation. Replacing
the obstacle matrix itself is a different operation: its minimizer can
have a different support. This section quantifies that distinction under
explicit sign and row-sum assumptions. All matrices and optima below are
comparison objects; no ambient preprocessing or local finder is supplied.

\begin{theorem}[Significant support is stable at squared-accuracy spectral error;
\statusproved]
\label{thm:op3-spectral-envelope-stability}
Let $0<e<1$, $\operatorname{vol}(G)>4/e$, $\lambda=e/2$, and
$b=e_v-\lambda d$. For $0\leq t\leq1$ put $M=(1-t)L+tD$.
Suppose $\widehat M$ is symmetric, has nonpositive off-diagonal entries
and nonnegative row sums, and, for $0<\eta\leq1/2$, satisfies
\[
 (1-\eta)M\preceq\widehat M\preceq(1+\eta)M.
\]
Let $u$ and $\widehat u$ be their nonnegative same-load obstacle minimizers.
Both exist uniquely, including at $t=0$. Their original support volumes
are at most $2/e$, and
\[
 \|u-\widehat u\|_\infty
 \leq\frac{\sqrt{32}\eta}{e(1-\eta)}<\frac{12\eta}{e}.
\]
Consequently $\eta\leq e^2/192$ guarantees
$\|u-\widehat u\|_\infty\leq e/16$ and
$\operatorname{supp}(\widehat u)\supseteq\{i:u_i>e/8\}$.
The degree vector and the support-volume measure are those of the original
unweighted graph, even when $\widehat M$ uses different conductances.
\end{theorem}
\begin{proof}
For $t>0$, both matrices are positive definite, giving unique orthant
minimizers. For $t=0$, spectral equivalence gives kernel exactly
$\operatorname{span}\{\one\}$ for both matrices; their nonpositive
off-diagonal entries make them connected weighted Laplacians. Moreover
$\one^Tb=1-\lambda\operatorname{vol}(G)<0$. Write a nonnegative
vector as $c\one+y$, with $\min y=0$. Its objective contains the
positive linear term $-c\one^Tb$ and a coercive quadratic-minus-linear
function of $y$. Coercivity follows because the Laplacian energy has
a positive minimum on the compact set $\{y\geq0:\min y=0,\max y=1\}$.
Existence follows. Two minimizers could differ only by a constant;
the nonzero constant-direction objective slope proves uniqueness.

For either minimizer $x$ and its corresponding matrix $B$, define
$r=e_v-Bx$. KKT gives $r_i=\lambda d_i$ on positive coordinates.
At a zero coordinate the nonpositive off-diagonal entries give $r_i\geq0$.
Also $\one^Tr=1-(B\one)^Tx\leq1$. Thus
$\lambda\operatorname{vol}(\operatorname{supp}x)\leq1$.
The union $W$ of the two supports therefore has volume at most $4/e$,
strictly below the original graph volume. It is proper, and $M_{WW}$
is positive definite even when $t=0$.

Write $z=\widehat u-u$ and $E=\widehat M-M$. The two variational
inequalities imply
\[
 z^TMz\leq-z^TE\widehat u
 \leq\eta\|z\|_M\|\widehat u\|_M.
\]
Here the spectral inequality bounds the bilinear form of $E$ in the
$M$-energy seminorm; equivalently restrict both vectors to $W$ and
use its positive definite principal matrix. The triangle inequality gives
\begin{equation}
 \|z\|_M\leq\frac{\eta}{1-\eta}\|u\|_M.
 \label{eq:op3-obstacle-spectral-energy-stability}
\end{equation}

We need bounds independent of a small $t$. For any vector $y$ supported
on $W$, telescope from $i\in W$ along a simple original path to the
first zero outside $W$. Cauchy--Schwarz gives
$y_i^2\leq |W|\,y^TLy$. If $t\leq1/2$, then $M\succeq L/2$,
so $y_i^2\leq2\operatorname{vol}(W)y^TMy$.
If $t\geq1/2$, then $M\succeq D/2$, giving the same coarse bound.
Thus, whenever $W$ is nonempty,
\[
 (M_{WW}^{-1})_{ii}\leq2\operatorname{vol}(W)\leq8/e,
 \qquad \|z\|_\infty^2\leq(8/e)\|z\|_M^2.
\]

For the original minimizer, $u^TMu=b^Tu\leq u_v\leq\|u\|_\infty$.
If $t\leq1/2$, the original unit-superlevel flux bound, also valid
at $t=0$, gives $\|u\|_\infty\leq
|\operatorname{supp}u|/(1-t)\leq2\operatorname{vol}(\operatorname{supp}u)$.
If $t\geq1/2$, the original mass identity gives $t\sum_i d_iu_i\leq1$,
so the same bound holds whenever the support is nonempty. Therefore
$\|u\|_M^2\leq4/e$. Combining this with
\cref{eq:op3-obstacle-spectral-energy-stability} proves the displayed
coordinate bound. If either needed union is empty, both vectors are zero
and the conclusion is immediate. The stated choice of $\eta$ yields
error at most $e/16$, which preserves strict positivity of every coordinate
originally above $e/8$.
\end{proof}

\begin{proposition}[The squared-accuracy power is necessary for this rule;
\statusproved]
\label{prop:op3-spectral-envelope-path-obstruction}
There is a sequence of original paths, accuracies $e\downarrow0$, positive
physical parameters $t=e^4/128$, and symmetric strictly diagonally dominant
M-matrices satisfying
\[
 M_t\preceq\widehat M_t\preceq(1+\eta_e)M_t,
 \qquad \widehat M_t\one=M_t\one=td,
 \qquad \eta_e/e^2\longrightarrow4,
\]
whose same-load obstacle $\widehat u_t$ omits a coordinate with
$u_{t,i}>3e/16$. In particular, the accuracy power two in
\cref{thm:op3-spectral-envelope-stability} is sharp up to constants
for direct matrix-replacement support containment. This is not a lower
bound on iterative refinement work or on a producer that emits extra labels.
\end{proposition}
\begin{proof}
Use the path and $\lambda,e,\delta$ from
\cref{thm:op3-fourth-power-support-obstruction}. Put $q=N^{-2}$ and
let $L_R$ be the unit Laplacian on just the edges
$(N,N+1),\ldots,(2N-1,2N)$. Set
\[
 \widehat L=L+qL_R,\qquad
 t=e^4/128,\qquad \widehat M_t=M_t+qL_R.
\]
The original load $b=e_N-\lambda d$ is unchanged.
We first exhibit the $\widehat L$ obstacle supported on
$T=\{1,\ldots,2N-1\}$ with zero endpoints. Edge $(i-1,i)$ has
conductance $c_i=1$ for $i\leq N$ and $c_i=1+q$ for $N<i\leq2N$.
Set
\[
 j_1=\frac{1-\lambda(4N-2+q(N-1))}{2+q},\qquad
 j_i=j_1+2\lambda(i-1)-\boldsymbol1_{i>N},\qquad
 z_i=\sum_{k=1}^i j_k/c_k.
\]
Summing the two arithmetic progressions gives $z_{2N}=0$.
The left fluxes are positive and the right fluxes negative, since
\[
 \frac1{8N^2+1}\leq j_1<\lambda,
 \qquad
 j_1-\lambda=\frac{1-2qN^2}{(2+q)(8N^2+1)}<0,
 \qquad j_{2N}<0.
\]
These inequalities follow by substitution for $N\geq4$.
Thus $z$ is positive precisely on $T$. Its second flux differences
give $\widehat Lz=b$ on $T$. At the omitted endpoint, the incoming
flux is $j_1<\lambda$. At $2N$ it is
$-j_{2N}=1-j_1-2\lambda(2N-1)\in(0,2\lambda]$.
All farther rows have zero flux. These are the full same-load KKT
conditions, using the original degrees one at $0$ and two at $2N$.

Since $\widehat M_tz=\widehat Lz+tAz\geq b$ and $z\geq0$,
the positive-parameter maximum principle gives $\widehat u_t\leq z$.
In particular $\widehat u_{t,0}=0$. For the original matrix, the safe
comparison in \cref{lem:op3-conservative-torsion-comparison} gives
$u_{t,0}>u_{0,0}-\delta/2=3\delta/2=3e/16$.

Finally, $0\preceq L_R\preceq L\preceq M_t/(1-t)$, so the
spectral claim holds with $\eta_e=q/(1-t)$; this tends to $4e^2$
in ratio as $N\to\infty$. For every $N\geq4$ it is below $5e^2$.
The extra Laplacian has zero row sums and nonpositive off-diagonal entries,
so both stated structural assumptions hold. Along this sequence, any
permitted relative tolerance asymptotically larger than $e^2$ allows
the displayed support omission.
\end{proof}

\begin{proposition}[Spectral equivalence does not imply the sign premise;
\statusproved]
\label{prop:op3-spectral-factor-sign-obstruction}
Arbitrarily accurate positive definite spectral approximations with
nonnegative row sums can have positive off-diagonal entries. Their
same-load obstacle KKT conditions need not imply nonnegative residuals
at zero coordinates. Hence the sign premise of
\cref{thm:op3-spectral-envelope-stability} is substantive.
\end{proposition}
\begin{proof}
On an original star with $m\geq3$ leaves, let $M=D-\gamma A$,
$0<\gamma<1$. Embed $C=\gamma^2(I-J/m)$ on the leaves, zero at the
center. Eliminating the center gives the exact decomposition
\[
 M=aa^T/m+(1-\gamma^2)I_{\rm leaves}+C,
 \qquad a=(m,-\gamma,\ldots,-\gamma)^T.
\]
Thus $0\preceq C\preceq M$. For any $0<\eta<1$, the matrix
$Z=M-\eta C$ satisfies $(1-\eta)M\preceq Z\preceq M$ and
$Z\one=M\one>0$. But every pair of distinct leaves has
$Z_{ij}=\eta\gamma^2/m>0$.

For example, take $\gamma=1/2$, $\lambda=3/4$, and a leaf seed $s$.
The vector with only $x_s=(1-\lambda)/Z_{ss}>0$ satisfies every
same-load KKT condition for $Z$. The center's inactive gradient is
$\lambda m-\gamma x_s>0$, and all other inactive gradients are positive.
At any other leaf $j$, however, the residual is
$r_j=-Z_{js}x_s<0$. This is an abstract spectral-guarantee
counterexample; no particular source algorithm is asserted to return $Z$.
\end{proof}

\paragraph{Source boundary and exact audit.}
The checked approximate-elimination source
\citep{kyng2016approximate} is compared in
\path{docs/literature/lcp-solvers.md}. Its factorization guarantee alone
does not supply the sign premise, a positivity-constrained pivot order,
or local discovery. The conclusions above permit spectral approximation
as a preconditioner for a sufficiently accurate computation.

\path{spectral_envelope_stability.py} checks 28,335 same-load pairs,
including one- and two-sided edge/grounding perturbations, 16,365 exact
union-Dirichlet diagonals, and 9,444 safe-threshold containments.
It checks 132 weighted paths through $N=4096$, 135,748 original-degree
load rows, three independent dense endpoint obstructions and 12 sign
counterexamples. The full run takes 39.683 seconds. These are proof
validators; perturbed-matrix residuals are not relabelled original ACL
certificates. Source/backend hashes and exact cases are in
\path{SPECTRAL_ENVELOPE_STABILITY_AUDIT.json}.
